% ============================================================================
% Fichier  : partie4/P04_C20_applications_cas.tex
% Titre    : Applications et Cas Intégrateurs
% Auteur   : MINKA MI NGUIDJOÏ Thierry Emmanuel
% Version  : 1.0 -- Juin 2026
% Statut   : RÉDIGÉ
% Sources  : M7_Cas_Pratiques_Integrateurs_PNC.docx (structure des cas)
%            Synthèse de l'ensemble du cours (Parties I à IV)
% Positionnement : Chapitre-capstone. Mobilise simultanément le PIU,
%                  le droit camerounais, les techniques d'acquisition,
%                  l'analyse système/réseau/mobile et la rédaction de rapport.
%                  Trois cas intégrateurs nouveaux + simulation d'examen.
%                  Pas de nouvelles connaissances : consolidation et
%                  application transversale.
% ============================================================================

\chapter{Applications et Cas Intégrateurs}
\label{chap:cas-int}

\epigraph{%
    « La théorie c'est quand on sait tout et que rien ne fonctionne.
    La pratique c'est quand tout fonctionne et que personne ne sait
    pourquoi. Ici, nous visons les deux. »%
}{--- Adapté d'Albert Einstein}

\niveauChapitre{Expert}{%
    Ce chapitre exige la maîtrise de l'ensemble des Parties~I à~IV.
    Il ne présente aucune notion nouvelle~; il consolide et
    intègre par la pratique.%
}

\begin{boxobjectifs}
\begin{itemize}
    \item Appliquer le PIU de bout en bout sur un cas complexe
          multi-sources et multi-juridictions.
    \item Mobiliser simultanément~: qualification juridique,
          techniques d'acquisition, analyse artefacts, corrélation
          temporelle et rédaction de rapport.
    \item Gérer les tensions CRO dans des situations de décision
          réelle sous contrainte de temps.
    \item Identifier les erreurs classiques d'investigation et
          les corriger.
    \item Se préparer à l'évaluation finale par une simulation
          complète de cas intégrateur.
\end{itemize}
\end{boxobjectifs}

\bigskip

% ============================================================================
\section*{Mode d'emploi de ce chapitre}
% ============================================================================

Ce chapitre fonctionne différemment des précédents~: il ne présente
pas de nouvelles techniques. Il vous place dans la position de
l'investigateur face à un cas complet, en mobilisant simultanément
tout ce que vous avez appris.

\textbf{Structure de chaque cas~:}
\begin{enumerate}
    \item \textbf{Brief initial}~: les faits tels qu'ils vous sont
          présentés (incomplets, parfois trompeurs~--- comme dans
          la réalité).
    \item \textbf{Phase d'investigation guidée}~: questions structurées
          mobilisant les différentes dimensions du cours.
    \item \textbf{Révélations progressives}~: des indices supplémentaires
          apparaissent à mesure que l'investigation avance.
    \item \textbf{Synthèse et rapport}~: production d'un livrable
          juridiquement structuré.
\end{enumerate}

\separateur

% ============================================================================
\section{Cas~1 --- Opération FAKO~: Fraude Comptable Interne}
\label{sec:cas-fako}
% ============================================================================

\subsection{Brief initial}

\begin{boxscenario}{Brief --- Opération FAKO}
\textbf{Date~:} Lundi 10~mars~2026, 09h15.

La DRH de CAMEROON AGRO~SA (exportateur de produits agricoles,
150~salariés, Buea) vous contacte en urgence. La direction a
découvert, lors de l'audit mensuel, que des virements bancaires
non autorisés totalisant \textbf{47~millions~FCFA} ont été effectués
vers des comptes inconnus entre le 1er et le 28~février~2026.

Informations disponibles~:
\begin{itemize}
    \item Les virements ont été initiés depuis le logiciel comptable
          interne (\texttt{Sage~100 Comptabilité}, serveur Windows~Server~2019)~;
    \item Trois comptes employés avaient accès au module de virements~:
          EBOA Claire (DAF), BAFUT Jonas (comptable senior),
          MBAH Inès (comptable junior)~;
    \item Aucune des trois n'a signalé d'anomalie~;
    \item MBAH Inès est en arrêt maladie depuis le 1er~mars~2026.
\end{itemize}

\textbf{Mission~:} vous êtes mandaté par le parquet suite à la plainte
de CAMEROON AGRO~SA. Identifier l'auteur des virements non autorisés.
\end{boxscenario}

\subsection{Phase~1~: Initialisation (PIU E1--E3)}

\begin{boxexo}[title={\ding{45}\quad FAKO --- Étape~1~: Cadrage}]
\begin{enumerate}[(a)]
    \item Remplissez la matrice de criticité (Tableau~6 du Ch.~PIU)
          pour cette affaire. Quel niveau de rigueur s'applique~?
    \item Quelles autorisations judiciaires devez-vous obtenir
          avant toute action technique~? Citez les articles exacts
          (\loicam{2010/012}).
    \item Listez les sources de données à collecter, dans l'ordre
          de la RFC~3227. Quelle est la première source à traiter~?
    \item Y a-t-il un risque de destruction de preuves actif~?
          Comment le gérer~?
\end{enumerate}
\end{boxexo}

\subsection{Révélation~1 --- Accès aux logs}

\begin{boiteRemarque}{Révélation~1 (disponible après l'Étape~1)}
L'administrateur système vous communique~: le serveur Sage~100 est
allumé~; les logs d'authentification Windows sont activés et conservés
sur 30~jours (Event~ID~4624/4634). La messagerie professionnelle
(Exchange) est hébergée sur Office~365 (Microsoft~365). MBAH Inès
avait accès au VPN depuis chez elle.
\end{boiteRemarque}

\subsection{Phase~2~: Acquisition et Analyse (PIU E4--E9)}

\begin{boxexo}[title={\ding{45}\quad FAKO --- Étape~2~: Investigation système}]
\begin{enumerate}[(a)]
    \item L'analyse des Event~ID~4624 (connexion réussie) sur le serveur
          Sage~100 révèle~:
\begin{lstlisting}[language=bash_forensique]
# Extrait Event Viewer -- Logon Events
[2026-02-14 02:17:43] EventID:4624 Account:MBAH_INES
                       LogonType:3 (Network)
                       SourceIP:41.202.XXX.XXX
[2026-02-14 02:17:45] EventID:4624 Account:MBAH_INES
                       LogonType:3 Source:41.202.XXX.XXX
[2026-02-21 03:05:12] EventID:4624 Account:MBAH_INES
                       LogonType:3 Source:41.202.XXX.XXX
\end{lstlisting}
          Que révèle l'heure des connexions (02h17 et 03h05)~?
          Quelle hypothèse formulez-vous~?
    \item L'IP \texttt{41.202.XXX.XXX} est localisée au Nigéria
          (Lagos). MBAH Inès est en arrêt maladie depuis le 1er~mars
          mais les connexions datent de février. Est-ce problématique~?
    \item Formulez 5~hypothèses (méthode ACH, E8 PIU). Quelle hypothèse
          est la plus cohérente avec les données actuelles~?
    \item Quelles données demandez-vous à Microsoft via le portail
          LEA pour corroborer ou réfuter~?
\end{enumerate}
\end{boxexo}

\subsection{Révélation~2 --- Résultat LEA Microsoft}

\begin{boiteRemarque}{Révélation~2}
La réponse de Microsoft (portail LEA, délai~: 5~jours) indique~:
\begin{itemize}
    \item Le compte \texttt{mbah.ines@cameroon-agro.cm} s'est connecté
          depuis l'IP \texttt{41.202.XXX.XXX} (Lagos, Nigéria) à
          02h17 le 14~février~; 1~email a été lu~: l'email de réinitialisation
          VPN envoyé par l'IT à MBAH Inès le 13~février.
    \item Le même compte s'est connecté depuis l'IP
          \texttt{196.168.XXX.XXX} (Douala, Cameroun) à 08h30
          le 14~février~--- 6~heures plus tard.
\end{itemize}
\end{boiteRemarque}

\subsection{Phase~3~: Qualification et rapport}

\begin{boxexo}[title={\ding{45}\quad FAKO --- Étape~3~: Qualification et rapport}]
\begin{enumerate}[(a)]
    \item La connexion depuis Lagos à 02h17 (accès email de réinitialisation
          VPN) suivie d'une connexion depuis Douala à 08h30 le même
          jour est-elle physiquement possible pour une personne
          se déplaçant entre les deux villes~? Quelle est l'implication
          forensique~?
    \item Qualifiez les infractions potentielles (articles Loi~2010/012
          et Loi~2024/017). Qui est susceptible d'être mis en cause~?
    \item Appliquez la grille CRO à la décision de notifier MBAH Inès
          qu'elle est suspectée avant l'obtention de son téléphone.
    \item Rédigez les 5~premières constatations de votre rapport
          (Section~3, Constatations), en citant chaque source.
\end{enumerate}
\end{boxexo}

% ============================================================================
\section{Cas~2 --- Opération MUNGO~: Trafic d'Influence en Ligne}
\label{sec:cas-mungo}
% ============================================================================

\subsection{Brief initial}

\begin{boxscenario}{Brief --- Opération MUNGO}
\textbf{Date~:} Mercredi 15~avril~2026, 14h00.

Le Service Central de Lutte contre la Cybercriminalité (SCFC) reçoit
un signalement~: une campagne coordonnée de désinformation cible un
candidat à des élections locales dans la région du Littoral. Des
centaines de faux comptes Facebook diffusent des \textit{deepfakes}
audio attribuant de fausses déclarations au candidat.

Informations disponibles~:
\begin{itemize}
    \item 347~faux comptes identifiés par les équipes de modération
          de Meta, actifs depuis le 1er~avril~2026~;
    \item Tous créés depuis des IPs dans un même sous-réseau
          (\texttt{196.200.XX.0/24}, hébergeur camerounais)~;
    \item Contenus~: 12~clips audio deepfake, 47~posts texte,
          3~faux articles de presse~;
    \item Plainte déposée par le candidat victime le 14~avril.
\end{itemize}
\end{boxscenario}

\subsection{Investigation guidée}

\begin{boxexo}[title={\ding{45}\quad MUNGO --- Investigation complète}]
\begin{enumerate}[(a)]
    \item \textbf{Cadrage juridique}~: qualifiez les infractions
          potentielles. Les deepfakes audio sont-ils couverts par
          la \loicam{2010/012}~? Par quelle disposition~?
          Quelle innovation de la \loicam{2024/017} est-elle pertinente
          (données biométriques dans les deepfakes)~?
    \item \textbf{Stratégie de conservation d'urgence}~: les faux
          comptes Facebook peuvent être supprimés par Meta en quelques
          heures (violation des conditions d'utilisation). Quelle est
          votre première action~? Rédigez la demande de conservation
          urgente (art.~29 Budapest, portail LEA Meta Emergency).
    \item \textbf{Analyse technique des deepfakes}~: les clips audio
          deepfake sont des fichiers~MP3. Quelle analyse forensique
          pouvez-vous conduire sur ces fichiers~? Quels outils~?
          Quels artefacts~? (Indice~: métadonnées EXIF audio,
          spectrogramme, détecteurs de deepfake IA).
    \item \textbf{Attribution des comptes}~: tous les comptes viennent
          du même sous-réseau \texttt{196.200.XX.0/24}. Comment
          identifier l'abonné de cet hébergeur camerounais~?
          Quelle réquisition~? Quels délais~?
    \item \textbf{Contexte électoral et urgence déontologique}~:
          l'élection a lieu dans 15~jours. Le responsable politique
          concurrent tente de vous contacter ``en privé'' pour
          ``accélérer le dossier''. Identifiez la tension déontologique
          et décrivez votre réponse (Ch.~\ref{chap:deonto}).
    \item \textbf{CRO}~: appliquez la grille CRO à la décision de
          publier l'analyse technique des deepfakes avant la fin
          de l'investigation~: qui en bénéficierait, qui en
          souffrirait~?
\end{enumerate}
\end{boxexo}

% ============================================================================
\section{Cas~3 --- Opération SANAGA~: Ransomware Hôpital}
\label{sec:cas-sanaga}
% ============================================================================

\subsection{Brief initial}

\begin{boxscenario}{Brief --- Opération SANAGA}
\textbf{Date~:} Samedi 22~mars~2026, 06h30.

L'Hôpital Général de Yaoundé subit une attaque ransomware.
Les systèmes informatiques du bloc chirurgical, de la pharmacie
et des dossiers patients sont paralysés. Des interventions
chirurgicales ont dû être reportées. Un message de rançon
exige 15~bitcoins ($\approx$~750~000~USD) dans les 48~heures.

Particularités critiques~:
\begin{itemize}
    \item Des données de santé de 12~000~patients camerounais
          ont potentiellement été exfiltrées~;
    \item Le serveur principal est sous Windows~Server~2019,
          allumé, avec le message de rançon visible à l'écran~;
    \item Deux vies sont en jeu~: un patient en soins intensifs
          dépend d'un équipement connecté au réseau compromis~;
    \item Le chef du service informatique affirme avoir ``fait
          une sauvegarde hier''~--- mais ne retrouve pas le support.
\end{itemize}
\end{boxscenario}

\subsection{Décisions critiques sous pression}

\begin{boxexo}[title={\ding{45}\quad SANAGA --- Décisions immédiate et médiate}]
\textbf{Vous arrivez sur les lieux à 07h00. Vous avez 10~minutes
pour prendre vos premières décisions.}

\begin{enumerate}[(a)]
    \item \textbf{Triage}~: le patient en soins intensifs est-il
          votre responsabilité en tant qu'investigateur~?
          Que faites-vous en priorité absolue~? Citez le cadre
          éthique (Ch.~\ref{chap:deonto}~: tension efficacité
          vs droits, ici tension investigation vs vie humaine).
    \item \textbf{Ordre de volatilité}~: le serveur est allumé
          avec le message de rançon visible. Quelle est votre
          séquence d'actions dans les 30~premières minutes~?
          (RFC~3227, Ch.~\ref{chap:meth-std}). Argumentez pourquoi
          la capture RAM précède la capture réseau dans ce contexte.
    \item \textbf{Données de santé}~: 12~000~dossiers patients
          ont potentiellement été exfiltrés. Quelles obligations
          s'appliquent à l'hôpital en vertu de la
          \loicam{2024/017}~? Quels articles~? Dans quels délais~?
          Qui doit en être informé~?
    \item \textbf{La rançon}~: le Directeur de l'hôpital vous
          demande si vous ``conseillez'' de payer la rançon.
          Que répondez-vous~? Justifiez en distinguant le rôle
          de l'investigateur, le rôle du décideur et les
          implications légales du paiement.
    \item \textbf{Investigation post-incident}~:
\begin{lstlisting}[language=bash_forensique]
# Volatility 3 -- windows.pstree (extrait)
4892  4888  WannaCryptor.exe   2026-03-22 04:17:22
  4888 4884  cmd.exe            2026-03-22 04:17:20
    4884 4880  mshta.exe        2026-03-22 04:17:15
      4880 4876  winword.exe    2026-03-22 04:16:58
\end{lstlisting}
          Analysez cette arborescence~: quelle chaîne d'infection
          suggère-t-elle~? Quelle vulnérabilité est probablement
          exploitée~? Rédigez la constatation correspondante.
    \item \textbf{Notification et coopération}~: le message de
          rançon contient une adresse Bitcoin. Tracez la démarche
          complète~: analyse blockchain, identification de l'exchange
          de sortie, demande de gel, juridictions impliquées.
\end{enumerate}
\end{boxexo}

% ============================================================================
\section{Grille d'Erreurs Classiques~: Diagnostiquer une Investigation}
\label{sec:erreurs-classiques}
% ============================================================================

\subsection{L'investigation TANGUE revisitée~: trouver les erreurs}

\begin{boxscenario}{Investigation viciée --- à corriger}
Voici le procès-verbal d'investigation de l'Opération TANGUE
(version non corrigée). Identifiez toutes les erreurs.

\textit{``Le 14~mars~2026 à 14h30, l'OPJ BITA s'est rendu au
domicile de BELLO Armand et a saisi son PC portable HP. L'OPJ
a allumé le PC pour constater les faits~; l'écran de veille
montrait un dossier nommé ``fraudes''. L'OPJ a photographié
l'écran et a éteint le PC. Le lendemain, l'OPJ a confié le PC
à l'analyste SONG pour analyse~; aucun document de transfer
n'a été rédigé. L'analyste a ouvert l'image forensique avec
Windows Explorer pour naviguer rapidement avant d'ouvrir Autopsy.
Le rapport indique~: ``l'accusé a dissimulé 127 fichiers de
preuves dans un dossier caché~; il est coupable de fraude
informatique selon art.~72.'' Les hash du disque original
sont conservés uniquement dans le fichier E01.''}}
\end{boxscenario}

\begin{boxexo}[title={\ding{45}\quad Diagnostic d'investigation viciée}]
\begin{enumerate}[(a)]
    \item Identifiez et nommez toutes les violations de procédure
          (au moins~7). Pour chacune, citez la règle violée
          (Ch.~\ref{chap:custody}, ACPO, ISO~27037).
    \item Lesquelles de ces violations rendent les preuves
          \textbf{irrecevables} devant un tribunal~? Justifiez.
    \item Lesquelles peuvent être \textbf{régularisées a posteriori}~?
          Comment~?
    \item Réécrivez le paragraphe du rapport en corrigeant les
          erreurs de rédaction (constatation vs analyse,
          qualification juridique par l'expert).
\end{enumerate}
\end{boxexo}

% ============================================================================
\section{Simulation d'Examen Intégrateur}
\label{sec:simulation-examen}
% ============================================================================

\begin{boxalert}
\textbf{Cette section constitue une évaluation formelle.}
Elle est conçue pour être utilisée en conditions d'examen (3~heures)
ou comme devoir maison complet. Toutes les réponses doivent citer
les articles de loi, les standards, les outils et les chapitres
du cours. Aucun document ni outil n'est autorisé en condition examen
sauf indication contraire du formateur.
\end{boxalert}

\begin{boxscenario}{Cas d'examen --- Opération DIBAMBA}
\textbf{Contexte~:} Le 5~mai~2026 à 23h45, le SCFC reçoit un
signalement de MTN Cameroun~: des transactions MoMo anormales
touchent 234~clients en 20~minutes. Montant total~: 68~millions~FCFA.

\textbf{Éléments techniques disponibles~:}
\begin{enumerate}
    \item Les transactions proviennent d'un seul numéro MTN
          (677~XXXXXX)~; le titulaire du numéro est EKINDI Paul,
          adresse~: Bonabéri, Douala~;
    \item Ce numéro a subi un SIM~swap à 23h30 (15~minutes avant
          les premières transactions)~;
    \item Le téléphone associé au numéro avant le swap était un
          Samsung Galaxy A32 (IMEI~: 358419XXXXXXXXX)~;
    \item Une capture réseau du réseau GSM (opération d'écoute
          légale, ordonnance N°~XXX) révèle des communications
          DATA depuis l'IMEI 358419XXXXXXXXX vers l'IP
          \texttt{41.66.XX.XX} (hébergeur OVH Paris) sur le
          port~8443 à 23h31~;
    \item L'adresse IP \texttt{41.66.XX.XX} héberge un domaine
          enregistré le 29~avril~2026~: \texttt{mtn-verification.cm}.
\end{enumerate}

\textbf{Contraintes~:}
\begin{itemize}
    \item Il est 01h00 du matin~; les opérateurs humains sont limités~;
    \item MTN Cameroun a conservé les logs de la transaction
          SIM~swap pendant 90~jours~;
    \item Le parquet de permanence peut délivrer une réquisition
          d'urgence~;
    \item Vous disposez de 2~heures avant que les logs de connexion
          de l'hébergeur OVH ne risquent d'être en rotation.
\end{itemize}
\end{boxscenario}

\begin{boxexo}[title={\ding{45}\quad Examen --- Opération DIBAMBA (3 heures)}]
\textbf{Partie~A --- Cadrage et urgences (45 minutes)}

\begin{enumerate}[(a)]
    \item Remplissez la matrice de criticité. Quel niveau de rigueur~?
    \item Listez 3~actions d'urgence (dans l'ordre) à réaliser
          dans les 30~premières minutes, avec base légale pour chacune.
    \item Rédigez la demande de conservation urgente à OVH (portail
          LEA, art.~29 Budapest / art.~57 \loicam{2010/012}).
          La France est partie à Budapest depuis 2006.
\end{enumerate}

\textbf{Partie~B --- Analyse technique (60 minutes)}

\begin{enumerate}[(a)]
    \setcounter{enumi}{3}
    \item Le SIM~swap s'est produit à 23h30. Comment cela permet-il
          à l'attaquant d'accéder aux comptes MoMo de 234~victimes~?
          Expliquez le mécanisme technique (OTP par SMS).
    \item La capture réseau montre une connexion vers le port~8443
          sur \texttt{mtn-verification.cm} (domaine enregistré
          le 29~avril). Que peut contenir cette connexion~?
          Quels filtres Wireshark utilisez-vous~?
    \item L'analyse OSINT du domaine \texttt{mtn-verification.cm}
          révèle~: registrar américain, paiement en Bitcoin,
          hébergement OVH Paris. Quelle séquence de réquisitions
          préparez-vous~?
\end{enumerate}

\textbf{Partie~C --- Qualification et rapport (75 minutes)}

\begin{enumerate}[(a)]
    \setcounter{enumi}{6}
    \item Qualifiez toutes les infractions potentielles~: pour chaque
          infraction, citez l'article exact (\loicam{2010/012},
          \loicam{2024/017}, Code pénal). Y a-t-il des infractions
          à la charge de MTN Cameroun pour la gestion du SIM~swap~?
    \item Formulez 5~constatations et 2~analyses (avec niveau de
          confiance) à partir des éléments disponibles.
    \item Rédigez la Synthèse Exécutive (Niveau~1) destinée
          au Procureur de permanence~: 1~page maximum,
          3~points, une pièce clé.
    \item \textbf{[Question bonus]} Appliquez la grille CRO à la
          décision de publier l'IP~\texttt{41.66.XX.XX} dans un
          communiqué de presse avant l'arrestation~: quelles
          tensions~? Quelle recommandation~?
\end{enumerate}
\end{boxexo}

\separateur

\begin{boxresume}
\textbf{Ce que ce chapitre vous a demandé de mobiliser simultanément~:}
\begin{itemize}
    \item \textbf{Droit camerounais}~: qualification précise des
          infractions (\loicam{2010/012}, \loicam{2024/017},
          Code pénal), règles d'articulation, sursis impossible
          (art.~89).

    \item \textbf{PIU}~: les phases E1--E17, les Quality Gates,
          la boucle itérative, la falsification active (E9bis),
          la distinction constatation/analyse.

    \item \textbf{Standards}~: RFC~3227 (ordre de volatilité),
          ISO~27037 (custody), ACPO~4~principes, NIST~SP~800-86
          (rapport).

    \item \textbf{Techniques}~: Event~ID Windows, logs SSH, Volatility
          (pstree, netscan), analyse portail LEA, conservation urgente
          art.~29 Budapest, analyse blockchain, corrélation temporelle.

    \item \textbf{Déontologie}~: tension efficacité vs vie humaine
          (SANAGA), pression politique (MUNGO), coverage d'un
          collègue (TANGUE).

    \item \textbf{Rapport}~: trois niveaux, 7~pièges, Annexe~G,
          formule de limitation de périmètre, qualification
          technique sans qualification juridique.
\end{itemize}

\textbf{Les trois leçons transversales de ce chapitre~:}
\begin{enumerate}
    \item \textbf{La conservation d'urgence prime toujours~:}
          dans FAKO, MUNGO, SANAGA et DIBAMBA, la première
          action est toujours de préserver avant que les données
          disparaissent. Aucune analyse ne compense des données
          perdues.
    \item \textbf{Les décisions techniques ont des conséquences
          déontologiques~:} le choix d'éteindre ou non un serveur
          (SANAGA) engage une responsabilité éthique au-delà de
          la technique.
    \item \textbf{Un rapport sans hash n'existe pas juridiquement~:}
          l'Annexe~G est la fondation de toute opposabilité.
          Sans elle, toutes vos analyses sont des opinions,
          pas des preuves.
\end{enumerate}
\end{boxresume}

\bigskip

\begin{boxinfo}
\textbf{Fin de la Partie~IV.} Les vingt chapitres couverts dans
ce cours (Parties~I à~IV) constituent le socle complet de la
forensique numérique au niveau master~: fondements philosophiques,
processus, standards, droit camerounais, techniques d'acquisition
et d'analyse (système, réseau, mobile, cloud, IoT), rédaction
de rapport et intégration par la pratique.

La \textbf{Partie~V} --- si elle est incluse dans votre programme
--- approfondit les sujets prospectifs~: OSINT avancé,
anti-forensique et contre-mesures, intelligence artificielle
appliquée à la forensique, et perspectives post-quantiques.
Ces chapitres s'inscrivent dans la continuité directe des
techniques maîtrisées ici.

\textit{Les affaires de ce cours~--- MBONGO, MVOM, WOURI,
ONAMBELE, BAOBAB, TANGUE, MOABI, ROUSSEAU, et les nouveaux
cas FAKO, MUNGO, SANAGA, DIBAMBA~--- sont fictives mais
construites sur des schémas d'infractions réels. Chaque
technique, chaque article de loi, chaque outil présenté est
authentique et applicable dès demain sur le terrain.}
\end{boxinfo}
